\section{Continuous Assessment 1 (CA-1) --- Model Paper Set B}
\LPUPaperHeader{Continuous Assessment 1 (CA-1)}{Set B}{60 Minutes}{30 Marks}{CO1 (DC Circuits), CO2 (AC Circuits)}

\begin{center}
\sffamily\textbf{\color{LPUOrange}Note: Answer any FIVE questions. Each question carries 6 marks.}
\end{center}

\questionitem{1}{
    Using Nodal Analysis, find the voltages at nodes 1 and 2 for a circuit with two current sources: $I_{s1} = 5\text{ A}$ injecting into node 1, $I_{s2} = 3\text{ A}$ leaving node 2, with conductances $G_1 = 0.5\text{ S}$, $G_2 = 0.25\text{ S}$, and mutual conductance $G_{12} = 0.2\text{ S}$.
}{6}{CO1}{L3 Apply}

\questionitem{2}{
    State Norton's theorem. A linear DC network has open-circuit voltage $V_{oc} = 48\text{ V}$ and short-circuit current $I_{sc} = 6\text{ A}$. Construct its Norton equivalent circuit and compute the current delivered to a $12\,\Omega$ load.
}{6}{CO1}{L3 Apply}

\questionitem{3}{
    Derive the resonant frequency, quality factor $Q$, and bandwidth for a series $RLC$ circuit. If $R = 4\,\Omega$, $L = 20\text{ mH}$, and $C = 5\,\mu\text{F}$, calculate $f_0$, $Q_0$, and the half-power bandwidth $BW$.
}{6}{CO2}{L3 Apply}

\questionitem{4}{
    Two impedances $\mathbf{Z}_1 = 6 + j8\,\Omega$ and $\mathbf{Z}_2 = 8 - j6\,\Omega$ are connected in parallel across a $200\text{ V}$, $50\text{ Hz}$ supply. Compute: (a) Branch currents $\mathbf{I}_1$ and $\mathbf{I}_2$, (b) Total supply current $\mathbf{I}_T$, (c) Overall power factor, and (d) Total active power.
}{6}{CO2}{L3 Apply}

\questionitem{5}{
    Explain the Two-Wattmeter method for power measurement in a 3-phase balanced system. Prove that $\tan\phi = \sqrt{3}\frac{W_1 - W_2}{W_1 + W_2}$ and state the readings when the power factor is $0.5$ lagging.
}{6}{CO2}{L4 Analyze}

\questionitem{6}{
    In a balanced delta-connected load with phase impedance $\mathbf{Z}_{ph} = 20 + j15\,\Omega$ connected across a $400\text{ V}$, 3-phase supply, determine: (a) Phase voltage and line voltage, (b) Phase current and line current, and (c) Total power consumed.
}{6}{CO2}{L3 Apply}

\newpage
\subsection*{Solutions and Detailed Marking Rubrics --- CA-1 Set B}

\begin{solutionbox}[Solution to Question 3]
\textbf{Step 1: Resonant Frequency Calculation}\\
\[ f_0 = \frac{1}{2\pi\sqrt{LC}} = \frac{1}{2\pi\sqrt{20 \times 10^{-3} \times 5 \times 10^{-6}}} = \frac{1}{2\pi\sqrt{10^{-7}}} = \frac{1}{2\pi \times 3.162 \times 10^{-4}} \approx 503.29\text{ Hz} \]
\textbf{Step 2: Quality Factor ($Q_0$)}\\
\[ Q_0 = \frac{1}{R}\sqrt{\frac{L}{C}} = \frac{1}{4}\sqrt{\frac{20 \times 10^{-3}}{5 \times 10^{-6}}} = \frac{1}{4}\sqrt{4000} = \frac{63.245}{4} \approx 15.81 \]
\textbf{Step 3: Half-Power Bandwidth ($BW$)}\\
\[ BW = \frac{f_0}{Q_0} = \frac{503.29}{15.81} \approx 31.83\text{ Hz} \quad \text{or} \quad BW = \frac{R}{2\pi L} = \frac{4}{2\pi \times 0.020} = 31.83\text{ Hz} \]
\begin{markingrubric}
$\bullet$ Mathematical formula derivation for $f_0$, $Q$, and $BW$: \textbf{3.0 Marks}\\
$\bullet$ Correct numerical substitution and unit accuracy: \textbf{3.0 Marks}
\end{markingrubric}
\end{solutionbox}
\newpage
